\section{Tighter Strang Bound via Norm-of-Difference}
\label{apd:tighter}

We show that the Strang commutator-scaling bound can be tightened by
delaying the triangle inequality.
The standard bound
(\cref{sec:strang}, following Childs et al.~\cite{childsTheoryTrotterError2021})
uses a \emph{sum of norms}:
\begin{equation}\label{eq:strang-sum}
  \norm{S_2(t) - e^{tH}}
    \le \frac{\norm{\comm{B}{\comm{B}{A}}}}{12}\,t^3
      + \frac{\norm{\comm{A}{\comm{A}{B}}}}{24}\,t^3.
\end{equation}
We prove a bound using the \emph{norm of a difference}, whose leading
coefficient is always $\le$ the sum-of-norms coefficient and strictly
smaller when the two double commutators partially cancel.
In the proportional model of \cref{eq:tighten-ratio} the cancellation
condition becomes positive alignment; in general no alignment is required,
and the numerical sweep of \cref{sec:comparison-standard} finds strictly
positive cancellation in every tested spin-chain configuration, including
at zero and negative alignment proxy.

\subsection{Setup}

Write $A' = A/2$, $H = A + B$, and recall the Strang splitting
$S_2(t) = e^{tA'}\,e^{tB}\,e^{tA'}$.
The Duhamel approach (\cref{sec:commutator}) gives
\begin{equation}\label{eq:strang-duhamel}
  S_2(t) - e^{tH} = e^{tH} \int_0^t \mathcal{T}_2(\tau)\,d\tau
  \qquad \text{(anti-Hermitian case: $\norm{e^{sX}} = 1$)},
\end{equation}
where $\mathcal{T}_2(\tau) = \text{Part}_A(\tau) + \text{Part}_B(\tau)$
is the Strang residual from the 4-factor product rule
(\cref{sec:strang}).

\subsection{Extracting the leading order}

The two parts have the following structure
(from the double FTC, \cref{sec:commutator}):

\paragraph{Part~A (from $A'$-conjugation).}
Define $g_1(s) = e^{sA'}\,\comm{A'}{B}\,e^{-sA'}$.
Then $\text{Part}_A(\tau) = \int_0^\tau (g_1(s) - g_1(\tau))\,ds$
(the ``subtract-constant-at-$\tau$'' trick).

Expanding $g_1(s) - g_1(\tau) = -\int_s^\tau g_1'(u)\,du$
where $g_1'(u) = e^{uA'}\,\comm{A'}{\comm{A'}{B}}\,e^{-uA'}$:
\begin{align}
  g_1(s) - g_1(\tau)
    &= -(\tau - s)\,\comm{A'}{\comm{A'}{B}} - \int_s^\tau
        \bigl(g_1'(u) - \comm{A'}{\comm{A'}{B}}\bigr)\,du
        \notag \\
    &= -(\tau - s)\,\comm{A'}{\comm{A'}{B}} + O(\tau^2).
        \label{eq:g1-expand}
\end{align}
The second term is bounded via the single FTC applied to $g_1'$:
\[
  g_1'(u) - \comm{A'}{\comm{A'}{B}}
    = \int_0^u e^{vA'}\,\comm{A'}{\comm{A'}{\comm{A'}{B}}}\,e^{-vA'}\,dv,
\]
so $\norm{g_1'(u) - \comm{A'}{\comm{A'}{B}}} \le
  \norm{\comm{A'}{\comm{A'}{\comm{A'}{B}}}}\,|u|$
(anti-Hermitian isometry).

Integrating \eqref{eq:g1-expand}:
\begin{equation}\label{eq:partA-expand}
  \text{Part}_A(\tau) = -\frac{\tau^2}{2}\,\comm{A'}{\comm{A'}{B}}
    + R_A(\tau),
\end{equation}
where $R_A(\tau)$ is a triple integral satisfying
\begin{equation}\label{eq:RA-bound}
  \norm{R_A(\tau)} \le
    \frac{\norm{\comm{A'}{\comm{A'}{\comm{A'}{B}}}}}{3}\,\tau^3.
\end{equation}
The bound \eqref{eq:RA-bound} follows from
$\norm{R_A(\tau)} \le \int_0^\tau \int_s^\tau |u|\,du\,ds
  \cdot \norm{\comm{A'}{\comm{A'}{\comm{A'}{B}}}}
  = \frac{\tau^3}{3}\,\norm{\comm{A'}{\comm{A'}{\comm{A'}{B}}}}$.

\paragraph{Part~B (from $B$-conjugation).}
The double-FTC remainder
$R_2(\tau) = e^{\tau B}\,A'\,e^{-\tau B} - A' - \tau\,\comm{B}{A'}$
satisfies (by the double-integral representation):
\begin{equation}\label{eq:R2-expand}
  R_2(\tau) = \frac{\tau^2}{2}\,\comm{B}{\comm{B}{A'}} + R_B(\tau),
  \qquad
  \norm{R_B(\tau)} \le
    \frac{\norm{\comm{B}{\comm{B}{\comm{B}{A'}}}}}{6}\,\tau^3.
\end{equation}
The bound on $R_B$ is exactly our triple-FTC result
(\texttt{norm\_exp\_conj\_sub\_comm2\_le\_of\_skewAdjoint}
in \leansrc{HigherCommutator.lean}).

Since $\text{Part}_B(\tau) = e^{\tau A'}\,R_2(\tau)\,e^{-\tau A'}$,
conjugation by $e^{\tau A'}$ introduces an additional correction:
\begin{equation}\label{eq:partB-expand}
  \text{Part}_B(\tau) = \frac{\tau^2}{2}\,\comm{B}{\comm{B}{A'}}
    + \underbrace{\frac{\tau^2}{2}\bigl(
      e^{\tau A'}\comm{B}{\comm{B}{A'}}e^{-\tau A'}
      - \comm{B}{\comm{B}{A'}}\bigr)}_{R_C(\tau)}
    + e^{\tau A'}\,R_B(\tau)\,e^{-\tau A'},
\end{equation}
where $R_C(\tau)$ is bounded by the single FTC on $\comm{B}{\comm{B}{A'}}$:
\begin{equation}\label{eq:RC-bound}
  \norm{R_C(\tau)} \le
    \frac{\norm{\comm{A'}{\comm{B}{\comm{B}{A'}}}}}{2}\,\tau^3.
\end{equation}

\subsection{The tighter bound}

Combining \eqref{eq:partA-expand} and \eqref{eq:partB-expand}:
\begin{equation}\label{eq:T2-refined}
  \mathcal{T}_2(\tau) = \frac{\tau^2}{2}\,D + R(\tau),
\end{equation}
where the \emph{effective double commutator} is
\begin{equation}\label{eq:D-def}
  \boxed{D \;=\; \comm{B}{\comm{B}{A'}} - \comm{A'}{\comm{A'}{B}}}
\end{equation}
and $R(\tau) = R_A(\tau) + R_C(\tau) + e^{\tau A'}\,R_B(\tau)\,e^{-\tau A'}$
satisfies
\begin{equation}\label{eq:R-bound}
  \norm{R(\tau)} \le T \cdot \tau^3,
  \qquad
  T = \frac{\norm{\comm{A'}{\comm{A'}{\comm{A'}{B}}}}}{3}
    + \frac{\norm{\comm{A'}{\comm{B}{\comm{B}{A'}}}}}{2}
    + \frac{\norm{\comm{B}{\comm{B}{\comm{B}{A'}}}}}{6}.
\end{equation}

Integrating:
\begin{align}
  \norm{S_2(t) - e^{tH}}
    &\le \int_0^t \norm{\mathcal{T}_2(\tau)}\,d\tau
    \le \int_0^t \Bigl(\frac{\norm{D}}{2}\,\tau^2 + T\,\tau^3\Bigr)\,d\tau
    \notag \\
    &= \frac{\norm{D}}{6}\,t^3 + \frac{T}{4}\,t^4.
    \label{eq:strang-tighter}
\end{align}

Substituting $A' = A/2$ (using $\comm{cX}{Y} = c\,\comm{X}{Y}$ for
central scalars, so $\comm{A'}{\comm{A'}{B}} = \tfrac{1}{4}\comm{A}{\comm{A}{B}}$
and $\comm{B}{\comm{B}{A'}} = \tfrac{1}{2}\comm{B}{\comm{B}{A}}$):

\begin{equation}\label{eq:strang-tighter-final}
  \boxed{
  \norm{S_2(t) - e^{tH}}
    \le \frac{1}{6}\left\|\frac{\comm{B}{\comm{B}{A}}}{2}
      - \frac{\comm{A}{\comm{A}{B}}}{4}\right\| t^3
    + \frac{T'}{4}\,t^4
  }
  \qquad (t \ge 0,\ A^* = -A,\ B^* = -B)
\end{equation}
where $T'$ is the triple-commutator sum \eqref{eq:R-bound} evaluated at
$A' = A/2$. As throughout this appendix, the bound is stated for skew-adjoint
$A, B$ in a C\textsuperscript{*}-algebra and for $t \ge 0$; these are exactly the
hypotheses of the Lean theorem.

\subsection{Comparison with the standard bound}
\label{sec:comparison-standard}

\paragraph{Relation.}
The standard bound \eqref{eq:strang-sum} can be written as
\[
  \frac{\norm{\comm{B}{\comm{B}{A}}}}{12}
  + \frac{\norm{\comm{A}{\comm{A}{B}}}}{24}
  = \frac{1}{6}\left(
    \frac{\norm{\comm{B}{\comm{B}{A}}}}{2}
    + \frac{\norm{\comm{A}{\comm{A}{B}}}}{4}
    \right).
\]
By the triangle inequality:
\[
  \left\|\frac{\comm{B}{\comm{B}{A}}}{2}
    - \frac{\comm{A}{\comm{A}{B}}}{4}\right\|
  \le
  \frac{\norm{\comm{B}{\comm{B}{A}}}}{2}
    + \frac{\norm{\comm{A}{\comm{A}{B}}}}{4},
\]
so the leading term of \eqref{eq:strang-tighter-final} is always
$\le$ the standard bound \eqref{eq:strang-sum}, and strictly smaller exactly
when the triangle inequality above is strict for the pair
$\bigl(\tfrac12\comm{B}{\comm{B}{A}},\, -\tfrac14\comm{A}{\comm{A}{B}}\bigr)$,
i.e.\ when the two double commutators partially cancel in $D$. How much
cancellation to expect is a property of the specific pair: the proportional
model at \cref{eq:tighten-ratio} below gives intuition, and the numerical
sweep at the end of this subsection measures it on standard spin chains.
This domination is itself
formalized: the displayed inequality is
\leanref{StrangCommutatorScalingTight.lean}{576}{norm\_D\_le\_sum\_of\_norms}
(axiom-free, no C*-structure), so the leading-coefficient domination is
machine-checked, not merely argued.

\paragraph{When is the tightening significant?}
The $O(t^4)$ correction in \eqref{eq:strang-tighter-final} is controlled
by triple commutator norms. Assuming $T' > 0$, for $t$ below the threshold
\[
  t^* = \frac{2}{3}\,
  \frac{\norm{\comm{B}{\comm{B}{A}}}/2 + \norm{\comm{A}{\comm{A}{B}}}/4
    - \norm{D}}{T'},
\]
the new bound \eqref{eq:strang-tighter-final} is tighter than
\eqref{eq:strang-sum}; if $T' = 0$ the correction term vanishes and the new
bound is at least as tight for every $t \ge 0$, and strictly tighter for
every $t > 0$ whenever
$\norm{D}$ lies strictly below the standard coefficient. In the Trotter
regime ($t = O(1/n)$ with
$n \to \infty$), the $O(t^4)$ correction is negligible and the tightening
from the leading term dominates.

\paragraph{How much tightening, and when.}
Write $X = \comm{B}{\comm{B}{A}}$ and $Y = \comm{A}{\comm{A}{B}}$, and suppose the
two double commutators are proportional, $Y \approx c\,X$ for a real $c$. Then
$D = \tfrac{X}{2} - \tfrac{Y}{4} = \bigl(\tfrac12 - \tfrac{c}{4}\bigr) X$, while the
standard bound's coefficient is $\tfrac{\norm{X}}{2} + \tfrac{\norm{Y}}{4}
= \bigl(\tfrac12 + \tfrac{\abs{c}}{4}\bigr)\norm{X}$. The tightening ratio is therefore
\begin{equation}\label{eq:tighten-ratio}
  r(c) \;=\;
  \frac{\norm{D}}{\norm{X}/2 + \norm{Y}/4}
  \;=\;
  \frac{\abs{2 - c}}{2 + \abs{c}} .
\end{equation}
Two consequences hold \emph{within this proportional model}. First,
$r(c) = 1$ for every $c \le 0$: proportional anti-alignment adds
constructively in the difference $D$, and the norm-of-difference bound then
gives no improvement. Second, the model's gain comes only from
\emph{positive} alignment, and it is total at $c = 2$ ---
i.e.\ $\comm{A}{\comm{A}{B}} \approx 2\,\comm{B}{\comm{B}{A}}$ makes $D = 0$ and
annihilates the entire $t^3$ term. A $37.5\%$ reduction ($r = 5/8$) occurs at
$c = 6/13 \approx 0.46$. The model is diagnostic, not a characterization:
generic pairs of double commutators are not scalar multiples of one another,
strictness of the triangle inequality does not require scalar alignment, and
the sweep below finds strictly positive gains even in instances whose
alignment proxy is zero or negative.

\paragraph{Numerical check on standard spin chains.}
A sweep over spin-$\tfrac12$ open chains---Heisenberg XXZ
($\Delta \in \{0.5, 1, 2\}$, even/odd bond split), transverse-field Ising
($h \in \{0.5, 1, 2\}$, both even/odd and ZZ/field splits), and
longitudinal-plus-transverse-field Ising (ZZ/field split), at
$L = 4, 6, 8, 10$ sites, spectral norms---finds a strictly positive gain in
\emph{every} one of the 40 configurations tested: $\norm{D}$ falls below the
standard coefficient
$\tfrac12\norm{\comm{B}{\comm{B}{A}}} + \tfrac14\norm{\comm{A}{\comm{A}{B}}}$
by $9$--$50\%$ (quoted for the better of the two operator orderings, both of
which are computed). Here $c_{\mathrm{eff}} =
\operatorname{Re}\langle X, Y\rangle_F / \langle X, X\rangle_F$, with
$X = \comm{B}{\comm{B}{A}}$ and $Y = \comm{A}{\comm{A}{B}}$, is the
least-squares coefficient of $Y$ on $X$ in the Frobenius inner
product---an unnormalized regression diagnostic, not a cosine; real pairs
are not scalar multiples of one another. Heisenberg chains
under the even/odd split show positive alignment
($c_{\mathrm{eff}} \approx +0.15$ to $+0.5$; gains
$39$--$50\%$). Across the tested configurations the gain is not
alignment-driven: in the
TFIM and LTFIM ZZ/field splits the two double commutators are exactly
Frobenius-orthogonal ($c_{\mathrm{eff}} = 0$, where the proportional model
predicts no gain), yet the measured gains are $24$--$33\%$, and
configurations with $c_{\mathrm{eff}} < 0$ retain $9$--$16\%$. Over the
tested sizes the gain decays slowly (Heisenberg XXX,
even/odd split: $50.0\%$, $42.7\%$, $41.3\%$, $38.6\%$ at $L = 4, 6, 8, 10$),
and the operator ordering matters (e.g.\ $15.5\%$ vs.\ $31.3\%$ for the TFIM
ZZ/field split at $h = 2$), so in practice one evaluates $\norm{D}$ for both
orderings and takes the better---cheap at benchmark sizes, though dense
spectral norms grow exponentially with $L$, so at scale one would instead
bound $\norm{D}$ through interaction-graph locality. Four representative
configurations were re-derived by an independent reimplementation (reversed
tensor-ordering convention, explicit-SVD norms), agreeing to relative error
below $10^{-9}$; scripts and raw
data accompany the repository
(\texttt{scripts/sweep\_strang\_alignment.py},
\texttt{claude/strang\_alignment\_sweep.csv}).

\subsection{Why this tightening is absent from Childs et al.}
\label{sec:why-absent}

The standard bound \eqref{eq:strang-sum} appears in
Childs et al.~\cite{childsTheoryTrotterError2021} (Prop.~16 in the arXiv
version, specialized to two summands) and
throughout the product formula literature. Several reasons explain why
the norm-of-difference formulation \eqref{eq:strang-tighter-final} has
not appeared previously.

\paragraph{Proof architecture.}
Childs et al.\ use a ``conjugation expansion'' (their Theorem~5) that
expands the error at each exponential interface independently, producing
separate terms for the $A$-$B$ and $B$-$A$ boundaries. The triangle
inequality on the resulting sum is the most direct bound within this
framework. The norm-of-difference requires an additional step:
recombining the interface terms, factoring out the common $\tau^2/2$
prefactor, and bounding the remainder with a \emph{triple}-commutator
estimate (\cref{eq:R-bound}). This extra step does not arise naturally
from the interface-by-interface analysis.

\paragraph{Framework compatibility.}
Childs et al.\ build a systematic framework around the
\emph{nested-commutator sum}
$\tilde\alpha_{\mathrm{comm}} =
  \sum \norm{\comm{H_{\gamma_{p+1}}}{\ldots\comm{H_{\gamma_2}}{H_{\gamma_1}}}}$,
which extends uniformly to any order~$p$ and any number of operators.
The effective commutator $D = \comm{B}{\comm{B}{A'}} - \comm{A'}{\comm{A'}{B}}$
is a \emph{difference} of two nested commutators---it does not fit into
the $\tilde\alpha_{\mathrm{comm}}$ framework as a single nested commutator.

\paragraph{Simplicity.}
The standard bound \eqref{eq:strang-sum} is a clean $t^3$ expression
with no correction term, yielding a simple Trotter-number formula
$r \sim (C\,t^3/\varepsilon)^{1/2}$ for resource estimation.
The tighter bound \eqref{eq:strang-tighter-final} introduces an $O(t^4)$
correction involving triple commutators, making the resource formula
more complex. For asymptotic scaling analysis, the simpler expression
is often preferred.

\paragraph{When the tightening matters.}
Despite these considerations, the leading coefficient of the
norm-of-difference bound is never larger, and the sweep of
\cref{sec:comparison-standard} realizes the improvement throughout the
tested spin-chain configurations: every one gains
$9$--$50\%$, including instances with zero or negative alignment proxy.
Given a specific
Hamiltonian $H = A + B$, the effective commutator $D$ is computable
directly at benchmark sizes (at scale, locality-based bounds replace the
dense computation), so one need not
guess: evaluate $\norm{D}$ for both operator orderings and compare. In the
Trotter regime ($t = O(1/n)$),
the $O(t^4)$ correction is subdominant, so the leading-term tightening
is what survives.
The formalization-driven insight here is that the Duhamel approach---by
producing the residual $\mathcal{T}_2(\tau)$ as a \emph{single}
algebraic expression---makes the norm-of-difference formulation natural,
whereas the interface-by-interface expansion does not.

\subsection{Extension to higher-order formulas}

The same principle extends to the fourth-order Suzuki formula $S_4$.
Childs et al.~\cite{childsTheoryTrotterError2021} (Prop.~J.1 in the arXiv
version) bound
the $S_4$ error as a sum of 8 separate 4-fold commutator norms:
\[
  \norm{S_4(t) - e^{tH}}
    \le \sum_{k=1}^{8} \alpha_k\,\norm{C_k}\,t^5,
\]
with numerical coefficients $\alpha_k \in [0.0046, 0.0284]$.
The norm-of-difference formulation would replace this with
\[
  \norm{S_4(t) - e^{tH}} \le \norm{R_5}\,t^5 + O(t^6),
  \qquad
  R_5 = \sum_{k=1}^{8} \beta_k(p)\,C_k,
\]
where $R_5$ is the single algebraic error expression of \cref{eq:log-s4} and
the $\beta_k(p)$ are its \emph{signed} BCH coefficients \eqref{eq:bch-gamma}.
A note on the tail: palindromic symmetry removes the $t^6$ term from the
\emph{logarithm} ($\log S_4(t) - tH = t^5 R_5 + O(t^7)$), but exponentiating
reintroduces one in the operator difference,
$S_4(t) - e^{tH} = t^5 R_5 + \tfrac{t^6}{2}(H R_5 + R_5 H) + O(t^7)$, so
$O(t^6)$ is what holds in a general normed algebra. For skew-adjoint $A, B$
the sharper $O(t^7)$ tail should be recoverable: $S_4(t)$ is then unitary,
$\Delta(t) = \log S_4(t) - tH$ is skew-adjoint for small $t$, and the
unitary Duhamel estimate $\norm{e^{tH+\Delta} - e^{tH}} \le \norm{\Delta}$
gives $\norm{S_4(t) - e^{tH}} \le \norm{R_5}\,t^5 + O(t^7)$; this local
skew-adjointness argument is not yet formalized or script-checked, so we
state the conservative tail.
Since $\norm{R_5} \le \sum_k \abs{\beta_k}\,\norm{C_k}
\le \sum_k \gamma_k\,\norm{C_k} \le \sum_k \alpha_k\,\norm{C_k}$
by the triangle inequality and \cref{sec:cas}, its leading $t^5$ coefficient
is never larger than that of the Level~3 bound or of the standard one, and
is strictly smaller whenever the 4-fold nested commutators $C_k$ partially
cancel.

A companion numerical test on the same spin-chain families as the sweep of
\cref{sec:comparison-standard} (60 rows; \texttt{scripts/test\_r5\_gain.py},
\texttt{claude/s4\_r5\_gain.csv}) measures precisely this gap:
$\norm{R_5}/\sum_k \gamma_k \norm{C_k}$ ranges from $0.14$ to $0.92$---a
further $8$--$86\%$ reduction beyond Level~3, with Heisenberg chains under
the even/odd split at $71$--$86\%$. Since all $\beta_k \ge 0$ at the Suzuki
point $p = 1/(4-4^{1/3})$, the gain is operator-level cancellation among
the $C_k$ up to sub-percent ceiling slack (the ratio against
$\sum_k \abs{\beta_k}\norm{C_k}$ differs by less than $0.001$). On the same instances the certified Level-3 sum
$\sum_k \gamma_k \norm{C_k}$ itself undercuts Childs et al.'s
$\sum_k \alpha_k \norm{C_k}$ by roughly $24$--$45\times$ (measured
$23.6$--$45.4\times$)---well beyond the worst-case
termwise guarantee---because the operators weight the coefficients
non-uniformly.

\subsection{Formalization status}

The triple-FTC building block (\cref{eq:R2-expand},
\leansrc{HigherCommutator.lean}) is fully formalized with 0
\texttt{sorry} axioms, and so is the tighter Strang bound
\eqref{eq:strang-tighter-final} itself:
\leanref{StrangCommutatorScalingTight.lean}{335}{norm\_strang\_comm\_scaling\_tight}
is sorry-free, and \texttt{\#print axioms} on it returns only
\texttt{[propext, Classical.choice, Quot.sound]}.
The \emph{comparison} with the standard bound is machine-checked as well:
\leanref{StrangCommutatorScalingTight.lean}{576}{norm\_D\_le\_sum\_of\_norms}
proves $\norm{D} \le \tfrac12\norm{\comm{B}{\comm{B}{A}}}
+ \tfrac14\norm{\comm{A}{\comm{A}{B}}}$ (equivalently
$\norm{D}/6 \le \norm{\comm{B}{\comm{B}{A}}}/12
+ \norm{\comm{A}{\comm{A}{B}}}/24$, \cref{sec:comparison-standard}), so the
tight bound's leading coefficient provably never exceeds the standard one; it
too is axiom-free and needs no C*-structure. What remains open is the
\emph{$S_4$ extension} of the norm-of-difference idea --- the $\norm{R_5}$-based
bound of \cref{sec:why-absent}: every formalized $S_4$ bound in this development
uses the sum $\sum_i \gamma_i \norm{C_i}$, never $\norm{R_5}$; the measured
$8$--$86\%$ gap reported in \cref{sec:why-absent} is what formalizing it
would recover.
